% Options for packages loaded elsewhere
\PassOptionsToPackage{unicode}{hyperref}
\PassOptionsToPackage{hyphens}{url}
\PassOptionsToPackage{dvipsnames,svgnames,x11names}{xcolor}
\documentclass[
  11pt,
]{article}
\usepackage{xcolor}
\usepackage[left=3cm,right=3in,top=2cm,bottom=2cm]{geometry}
\usepackage{amsmath,amssymb}
\setcounter{secnumdepth}{5}
\usepackage{iftex}
\ifPDFTeX
  \usepackage[T1]{fontenc}
  \usepackage[utf8]{inputenc}
  \usepackage{textcomp} % provide euro and other symbols
\else % if luatex or xetex
  \usepackage{unicode-math} % this also loads fontspec
  \defaultfontfeatures{Scale=MatchLowercase}
  \defaultfontfeatures[\rmfamily]{Ligatures=TeX,Scale=1}
\fi
\usepackage{lmodern}
\ifPDFTeX\else
  % xetex/luatex font selection
\fi
% Use upquote if available, for straight quotes in verbatim environments
\IfFileExists{upquote.sty}{\usepackage{upquote}}{}
\IfFileExists{microtype.sty}{% use microtype if available
  \usepackage[]{microtype}
  \UseMicrotypeSet[protrusion]{basicmath} % disable protrusion for tt fonts
}{}
\makeatletter
\@ifundefined{KOMAClassName}{% if non-KOMA class
  \IfFileExists{parskip.sty}{%
    \usepackage{parskip}
  }{% else
    \setlength{\parindent}{0pt}
    \setlength{\parskip}{6pt plus 2pt minus 1pt}}
}{% if KOMA class
  \KOMAoptions{parskip=half}}
\makeatother
\usepackage{color}
\usepackage{fancyvrb}
\newcommand{\VerbBar}{|}
\newcommand{\VERB}{\Verb[commandchars=\\\{\}]}
\DefineVerbatimEnvironment{Highlighting}{Verbatim}{commandchars=\\\{\}}
% Add ',fontsize=\small' for more characters per line
\usepackage{framed}
\definecolor{shadecolor}{RGB}{248,248,248}
\newenvironment{Shaded}{\begin{snugshade}}{\end{snugshade}}
\newcommand{\AlertTok}[1]{\textcolor[rgb]{0.94,0.16,0.16}{#1}}
\newcommand{\AnnotationTok}[1]{\textcolor[rgb]{0.56,0.35,0.01}{\textbf{\textit{#1}}}}
\newcommand{\AttributeTok}[1]{\textcolor[rgb]{0.13,0.29,0.53}{#1}}
\newcommand{\BaseNTok}[1]{\textcolor[rgb]{0.00,0.00,0.81}{#1}}
\newcommand{\BuiltInTok}[1]{#1}
\newcommand{\CharTok}[1]{\textcolor[rgb]{0.31,0.60,0.02}{#1}}
\newcommand{\CommentTok}[1]{\textcolor[rgb]{0.56,0.35,0.01}{\textit{#1}}}
\newcommand{\CommentVarTok}[1]{\textcolor[rgb]{0.56,0.35,0.01}{\textbf{\textit{#1}}}}
\newcommand{\ConstantTok}[1]{\textcolor[rgb]{0.56,0.35,0.01}{#1}}
\newcommand{\ControlFlowTok}[1]{\textcolor[rgb]{0.13,0.29,0.53}{\textbf{#1}}}
\newcommand{\DataTypeTok}[1]{\textcolor[rgb]{0.13,0.29,0.53}{#1}}
\newcommand{\DecValTok}[1]{\textcolor[rgb]{0.00,0.00,0.81}{#1}}
\newcommand{\DocumentationTok}[1]{\textcolor[rgb]{0.56,0.35,0.01}{\textbf{\textit{#1}}}}
\newcommand{\ErrorTok}[1]{\textcolor[rgb]{0.64,0.00,0.00}{\textbf{#1}}}
\newcommand{\ExtensionTok}[1]{#1}
\newcommand{\FloatTok}[1]{\textcolor[rgb]{0.00,0.00,0.81}{#1}}
\newcommand{\FunctionTok}[1]{\textcolor[rgb]{0.13,0.29,0.53}{\textbf{#1}}}
\newcommand{\ImportTok}[1]{#1}
\newcommand{\InformationTok}[1]{\textcolor[rgb]{0.56,0.35,0.01}{\textbf{\textit{#1}}}}
\newcommand{\KeywordTok}[1]{\textcolor[rgb]{0.13,0.29,0.53}{\textbf{#1}}}
\newcommand{\NormalTok}[1]{#1}
\newcommand{\OperatorTok}[1]{\textcolor[rgb]{0.81,0.36,0.00}{\textbf{#1}}}
\newcommand{\OtherTok}[1]{\textcolor[rgb]{0.56,0.35,0.01}{#1}}
\newcommand{\PreprocessorTok}[1]{\textcolor[rgb]{0.56,0.35,0.01}{\textit{#1}}}
\newcommand{\RegionMarkerTok}[1]{#1}
\newcommand{\SpecialCharTok}[1]{\textcolor[rgb]{0.81,0.36,0.00}{\textbf{#1}}}
\newcommand{\SpecialStringTok}[1]{\textcolor[rgb]{0.31,0.60,0.02}{#1}}
\newcommand{\StringTok}[1]{\textcolor[rgb]{0.31,0.60,0.02}{#1}}
\newcommand{\VariableTok}[1]{\textcolor[rgb]{0.00,0.00,0.00}{#1}}
\newcommand{\VerbatimStringTok}[1]{\textcolor[rgb]{0.31,0.60,0.02}{#1}}
\newcommand{\WarningTok}[1]{\textcolor[rgb]{0.56,0.35,0.01}{\textbf{\textit{#1}}}}
\usepackage{graphicx}
\makeatletter
\newsavebox\pandoc@box
\newcommand*\pandocbounded[1]{% scales image to fit in text height/width
  \sbox\pandoc@box{#1}%
  \Gscale@div\@tempa{\textheight}{\dimexpr\ht\pandoc@box+\dp\pandoc@box\relax}%
  \Gscale@div\@tempb{\linewidth}{\wd\pandoc@box}%
  \ifdim\@tempb\p@<\@tempa\p@\let\@tempa\@tempb\fi% select the smaller of both
  \ifdim\@tempa\p@<\p@\scalebox{\@tempa}{\usebox\pandoc@box}%
  \else\usebox{\pandoc@box}%
  \fi%
}
% Set default figure placement to htbp
\def\fps@figure{htbp}
\makeatother
\setlength{\emergencystretch}{3em} % prevent overfull lines
\providecommand{\tightlist}{%
  \setlength{\itemsep}{0pt}\setlength{\parskip}{0pt}}
\usepackage{placeins, graphicx, geometry, longtable, fancyhdr, currfile, wrapfig,
multicol, caption, capt-of, lipsum, booktabs, pgfpages, titling, fancyhdr, currfile}
\usepackage[export]{adjustbox}
\usepackage{xcolor}
\usepackage[normalem]{ulem}
\pagestyle{fancy}
\pretitle{
\begin{flushright}
% \write18{wget https://github.com/rgt47/kickstart/blob/master/sudoku_apple.pdf}
\includegraphics[width=3cm,valign=c]{sudoku_apple.pdf}\\
\end{flushright}
\begin{flushleft} \LARGE }
\posttitle{\par\end{flushleft}
\vskip 0.5em}
\predate{\begin{flushleft}\large}
	\postdate{\par\end{flushleft}}
	\preauthor{\begin{flushleft}\large}
	\postauthor{\par\end{flushleft}}
\fancyfoot[L]{\currfilename} %put date in header
\fancyfoot[R]{\includegraphics[width=.8cm]{sudoku_apple.pdf}}
\fancyhead[L]{\today} %put current file in footer
\setlength{\headheight}{13.59999pt}

\newcommand\zztab[3]{
\noindent
\begin{center} 
\includegraphics[keepaspectratio,width=#2]{#1}
\end{center} 
\captionof{table}{#3}
}

\newcommand\zzfig[3]{
\noindent
\begin{center} 
\includegraphics[keepaspectratio, width=#2]{#1}
\end{center} 
\captionof{figure}{#3}
}
\rhead{Multiple engines}
\usepackage{bookmark}
\IfFileExists{xurl.sty}{\usepackage{xurl}}{} % add URL line breaks if available
\urlstyle{same}
\hypersetup{
  pdftitle={Multiple engine plotting Penguins},
  pdfauthor={R.G. Thomas},
  colorlinks=true,
  linkcolor={Maroon},
  filecolor={Maroon},
  citecolor={Blue},
  urlcolor={blue},
  pdfcreator={LaTeX via pandoc}}

\title{Multiple engine plotting Penguins}
\author{R.G. Thomas}
\date{}

\begin{document}
\maketitle

{
\hypersetup{linkcolor=}
\setcounter{tocdepth}{2}
\tableofcontents
}
\subsection{R Visualization}\label{r-visualization}

\begin{Shaded}
\begin{Highlighting}[]
\CommentTok{\# Create scatter plot with R}
\FunctionTok{ggplot}\NormalTok{(penguins\_clean, }\FunctionTok{aes}\NormalTok{(}\AttributeTok{x =}\NormalTok{ body\_mass\_g, }\AttributeTok{y =}\NormalTok{ bill\_depth\_mm, }\AttributeTok{color =}\NormalTok{ species)) }\SpecialCharTok{+}
  \FunctionTok{geom\_point}\NormalTok{(}\AttributeTok{alpha =} \FloatTok{0.8}\NormalTok{, }\AttributeTok{size =} \DecValTok{3}\NormalTok{) }\SpecialCharTok{+}
  \FunctionTok{scale\_color\_manual}\NormalTok{(}\AttributeTok{values =} \FunctionTok{c}\NormalTok{(}\StringTok{"Adelie"} \OtherTok{=} \StringTok{"darkorange"}\NormalTok{, }\StringTok{"Chinstrap"} \OtherTok{=} \StringTok{"purple"}\NormalTok{, }\StringTok{"Gentoo"} \OtherTok{=} \StringTok{"\#008080"}\NormalTok{)) }\SpecialCharTok{+}
  \FunctionTok{labs}\NormalTok{(}
    \AttributeTok{title =} \StringTok{"Bill Depth vs Body Mass by Penguin Species (R)"}\NormalTok{,}
    \AttributeTok{x =} \StringTok{"Body Mass (g)"}\NormalTok{,}
    \AttributeTok{y =} \StringTok{"Bill Depth (mm)"}\NormalTok{,}
    \AttributeTok{color =} \StringTok{"Species"}
\NormalTok{  ) }\SpecialCharTok{+}
  \FunctionTok{theme\_minimal}\NormalTok{()}
\end{Highlighting}
\end{Shaded}

\pandocbounded{\includegraphics[keepaspectratio]{multiple_engine_files/figure-latex/r-plot-1.pdf}}

\subsection{Python Visualization}\label{python-visualization}

\begin{Shaded}
\begin{Highlighting}[]
\CommentTok{\# Import libraries}
\ImportTok{import}\NormalTok{ pandas }\ImportTok{as}\NormalTok{ pd}
\ImportTok{import}\NormalTok{ matplotlib.pyplot }\ImportTok{as}\NormalTok{ plt}
\ImportTok{import}\NormalTok{ seaborn }\ImportTok{as}\NormalTok{ sns}

\CommentTok{\# Access the penguins data from R}
\NormalTok{penguins }\OperatorTok{=}\NormalTok{ r.py\_penguins}

\CommentTok{\# Create scatter plot with Python}
\NormalTok{plt.figure(figsize}\OperatorTok{=}\NormalTok{(}\DecValTok{8}\NormalTok{, }\DecValTok{5}\NormalTok{))}
\NormalTok{sns.scatterplot(}
\NormalTok{    data}\OperatorTok{=}\NormalTok{penguins,}
\NormalTok{    x}\OperatorTok{=}\StringTok{"body\_mass\_g"}\NormalTok{,}
\NormalTok{    y}\OperatorTok{=}\StringTok{"bill\_depth\_mm"}\NormalTok{,}
\NormalTok{    hue}\OperatorTok{=}\StringTok{"species"}\NormalTok{,}
\NormalTok{    palette}\OperatorTok{=}\NormalTok{\{}\StringTok{"Adelie"}\NormalTok{: }\StringTok{"darkorange"}\NormalTok{, }\StringTok{"Chinstrap"}\NormalTok{: }\StringTok{"purple"}\NormalTok{, }\StringTok{"Gentoo"}\NormalTok{: }\StringTok{"\#008080"}\NormalTok{\},}
\NormalTok{    alpha}\OperatorTok{=}\FloatTok{0.8}\NormalTok{,}
\NormalTok{    s}\OperatorTok{=}\DecValTok{80}
\NormalTok{)}

\NormalTok{plt.title(}\StringTok{"Bill Depth vs Body Mass by Penguin Species (Python)"}\NormalTok{)}
\NormalTok{plt.xlabel(}\StringTok{"Body Mass (g)"}\NormalTok{)}
\NormalTok{plt.ylabel(}\StringTok{"Bill Depth (mm)"}\NormalTok{)}
\NormalTok{plt.tight\_layout()}
\NormalTok{plt.show()}
\end{Highlighting}
\end{Shaded}

\pandocbounded{\includegraphics[keepaspectratio]{multiple_engine_files/figure-latex/python-analysis-1.pdf}}

\subsection{Julia Visualization}\label{julia-visualization}

\begin{Shaded}
\begin{Highlighting}[]
\ControlFlowTok{try}
    \CommentTok{\# Basic Julia plot {-} only runs if Julia is available}
    \ImportTok{using} \BuiltInTok{CSV}
    \ImportTok{using} \BuiltInTok{DataFrames}
    \ImportTok{using} \BuiltInTok{Plots}
    
    \CommentTok{\# Force the GR backend and set output format}
    \ConstantTok{ENV}\NormalTok{[}\StringTok{"GKSwstype"}\NormalTok{] }\OperatorTok{=} \StringTok{"100"}  \CommentTok{\# Ensure headless plotting works}
\NormalTok{    Plots.}\FunctionTok{gr}\NormalTok{()  }\CommentTok{\# Explicitly use GR backend}
    
    \CommentTok{\# Check if file exists}
    \ControlFlowTok{if} \FunctionTok{isfile}\NormalTok{(}\StringTok{"penguins\_clean.csv"}\NormalTok{)}
        \CommentTok{\# Read data and create scatter plot}
\NormalTok{        penguins }\OperatorTok{=}\NormalTok{ CSV.}\FunctionTok{read}\NormalTok{(}\StringTok{"penguins\_clean.csv"}\NormalTok{, DataFrame)}
        
        \CommentTok{\# Convert species to string type}
\NormalTok{        penguins.species }\OperatorTok{=} \FunctionTok{string}\NormalTok{.(penguins.species)}
        
        \CommentTok{\# Create scatter plot with hex color codes and store in variable}
\NormalTok{        plt }\OperatorTok{=} \FunctionTok{scatter}\NormalTok{(}
\NormalTok{            penguins.body\_mass\_g,}
\NormalTok{            penguins.bill\_depth\_mm,}
\NormalTok{            group}\OperatorTok{=}\NormalTok{penguins.species,}
\NormalTok{            markersize}\OperatorTok{=}\FloatTok{6}\NormalTok{,}
\NormalTok{            alpha}\OperatorTok{=}\FloatTok{0.8}\NormalTok{,}
\NormalTok{            markerstrokewidth}\OperatorTok{=}\FloatTok{0}\NormalTok{,}
\NormalTok{            color}\OperatorTok{=}\NormalTok{[}\StringTok{"\#FFA500"}\NormalTok{, }\StringTok{"\#800080"}\NormalTok{, }\StringTok{"\#008080"}\NormalTok{],  }\CommentTok{\# hex codes for orange, purple, teal}
\NormalTok{            xlabel}\OperatorTok{=}\StringTok{"Body Mass (g)"}\NormalTok{,}
\NormalTok{            ylabel}\OperatorTok{=}\StringTok{"Bill Depth (mm)"}\NormalTok{,}
\NormalTok{            title}\OperatorTok{=}\StringTok{"Bill Depth vs Body Mass by Penguin Species (Julia)"}\NormalTok{,}
\NormalTok{            legend}\OperatorTok{=:}\NormalTok{bottomright,}
\NormalTok{            size}\OperatorTok{=}\NormalTok{(}\FloatTok{800}\NormalTok{, }\FloatTok{500}\NormalTok{),}
\NormalTok{            dpi}\OperatorTok{=}\FloatTok{300}\NormalTok{,}
\NormalTok{            fmt}\OperatorTok{=:}\NormalTok{png  }\CommentTok{\# Force PNG output format}
\NormalTok{        )}
        
        \CommentTok{\# Explicitly display the plot}
        \FunctionTok{display}\NormalTok{(plt)}
        
        \CommentTok{\# Return the plot as the last expression}
\NormalTok{        plt}
    \ControlFlowTok{else}
        \FunctionTok{println}\NormalTok{(}\StringTok{"Error: penguins\_clean.csv file not found"}\NormalTok{)}
        \CommentTok{\# Create a simple error plot}
        \FunctionTok{plot}\NormalTok{(title}\OperatorTok{=}\StringTok{"Error: Data file not found"}\NormalTok{, legend}\OperatorTok{=}\ConstantTok{false}\NormalTok{)}
    \ControlFlowTok{end}
\ControlFlowTok{catch} \ConstantTok{e}
    \FunctionTok{println}\NormalTok{(}\StringTok{"Julia error: "}\NormalTok{, }\ConstantTok{e}\NormalTok{)}
    \CommentTok{\# Return something that can be displayed}
    \FunctionTok{plot}\NormalTok{(title}\OperatorTok{=}\StringTok{"Julia Error Occurred"}\NormalTok{, legend}\OperatorTok{=}\ConstantTok{false}\NormalTok{)}
\ControlFlowTok{end}
\end{Highlighting}
\end{Shaded}

Plot\{Plots.GRBackend() n=3\}

\end{document}
